\documentclass[conference]{IEEEtran}

\usepackage{times}
\usepackage{amsmath}
\usepackage{booktabs}
\usepackage{hyperref}
\usepackage{xspace}

\hyphenation{op-tical net-works semi-conduc-tor}

\newcommand{\nomic}{\texttt{nomic-embed-text-v1.5}\xspace}
\newcommand{\HR}[1]{\mathrm{HR}@{#1}}
\newcommand{\MRR}[1]{\mathrm{MRR}@{#1}}
\newcommand{\NDCG}[1]{\mathrm{NDCG}@{#1}}

\begin{document}

\title{Evaluating Hybrid Sparse-Dense Retrieval\\
for EU Climate Legal Documents}

\author{
  \IEEEauthorblockN{Rameez Wani}
  \IEEEauthorblockA{School of Computing\\
    Dublin City University\\
    Dublin, Ireland\\
    rameezshafat.wani2@mail.dcu.ie}
  \and
  \IEEEauthorblockN{Ananya Warior}
  \IEEEauthorblockA{School of Computing\\
    Dublin City University\\
    Dublin, Ireland\\
    ananyasalil.warior2@mail.dcu.ie}
}

\maketitle

\begin{abstract}
Legal information retrieval in statutory corpora poses specific
challenges due to formulaic language, hierarchical authority structures,
and the need for doctrinal precision.
This paper evaluates three first-stage retrieval configurations---BM25
lexical search, \nomic dense retrieval (768-dim,
8{,}192-token context), and a hybrid combining both via weighted
Reciprocal Rank Fusion (RRF, $k{=}20$, dense weight~5)---on a
manually constructed benchmark of 71~queries over EU Climate Law.
The corpus comprises 1{,}166 article-level provisions extracted from
72~EU legislative instruments via SPARQL and Formex~XML.
Evaluation uses a stratified 70/30 validation/test split; all reported
metrics are from the sealed 53-query test set.
Dense retrieval achieves $\HR{5}{=}0.925$, substantially outperforming
BM25 ($\HR{5}{=}0.698$; Wilcoxon $p{=}0.004$).
The hybrid achieves $\HR{5}{=}0.906$, $\MRR{5}{=}0.773$, and
$\NDCG{5}{=}0.721$, the highest MRR and NDCG of all three systems,
but its improvement over dense alone is statistically non-significant
($p{=}0.614$).
Only 5 of 53 queries ($9.4\,\%$) are hard failures for the hybrid.
The results confirm that a strong general-domain dense retriever
already approaches ceiling performance on this corpus, and that
weighted lexical-semantic fusion provides consistent ranking gains
without statistically significant coverage improvement.
\end{abstract}

\begin{IEEEkeywords}
Legal Information Retrieval, Hybrid Retrieval, BM25,
Dense Retrieval, Reciprocal Rank Fusion, EU Climate Law,
Retrieval-Augmented Generation, Statute-Level Benchmark
\end{IEEEkeywords}

\IEEEpeerreviewmaketitle

% ============================================================
\section{Introduction}

Legal information retrieval differs fundamentally from open-domain
search: the meaning of law is contingent on exact wording,
cross-references, and authority relationships.
In EU climate and energy legislation---which spans 72~enacted
instruments in the CELLAR repository with dense terminological
overlap---retrieval errors produce doctrinal misinterpretation rather
than mere semantic approximation.
This problem is consequential for retrieval-augmented generation (RAG)
pipelines, where a first-stage retriever that fails to surface the
correct statutory provision makes grounded, citation-supported
generation impossible~\cite{14}.

A recurring tension in legal IR research is the precision-abstraction
trade-off between lexical and semantic retrieval.
Lexical models such as BM25 are robust in formulaic corpora where
exact term overlap is indicative of relevance~\cite{1}, but fail when
user query vocabulary diverges from statutory language.
Dense models generalise beyond exact matches but introduce their own
failure modes: Document-Level Retrieval Mismatch (DRM), where
semantically similar passages from irrelevant instruments are retrieved
due to shared structural boilerplate~\cite{12}, and degraded
performance on specialised corpora without domain adaptation~\cite{14,16}.
Hybrid fusion addresses both limitations simultaneously~\cite{3}.

Prior work on hybrid legal retrieval has not systematically evaluated
statute-level EU climate law, which presents particular challenges:
cross-instrument definitional dependencies (e.g., CBAM cross-references
ETS allowance definitions from Directive~2003/87/EC), highly technical
vocabulary (\textit{free allocation}, \textit{tonne-kilometre data},
\textit{carbon leakage}), and a corpus where amendment instruments
contain Article~1 boilerplate that matches many queries weakly.
This paper addresses that gap.

\textbf{RQ1.} Does hybrid sparse-dense retrieval improve effectiveness
and reliability over BM25-only and dense-only baselines for EU Climate
Law statute retrieval?
\textbf{RQ2.} Is first-stage hybrid retrieval sufficient to support
grounded legal AI, given coverage, hard-failure rates, and ranking
quality?

Our contributions are: (i)~a 71-query manually verified benchmark for
EU Climate Law statute-level retrieval with binary CELEX-level
relevance judgements and a sealed 53-query test set; (ii)~a
reproducible 1{,}166-provision corpus from 72~EU climate instruments
extracted via a documented two-phase SPARQL pipeline; and (iii)~a
three-way comparative evaluation across effectiveness, reliability, and
statistical significance.

% ============================================================
\section{Related Work}

\subsection{Lexical vs.\ Dense Retrieval in Legal IR}

Mori et al.~\cite{1} demonstrate that BM25 remains a competitive
baseline for formulaic legal texts where exact term overlap is a
reliable relevance signal.
Dense models tend to underperform in such settings because semantic
similarity ignores legally binding phrasing.
Conversely, Zheng et al.~\cite{8} show that lexical models degrade on
reasoning-intensive queries; semantic models handle analogical
reasoning more effectively.
The LRAGE evaluation framework~\cite{14} confirms that off-the-shelf
dense retrievers underperform BM25 on legal retrieval tasks without
domain-specific adaptation.
Li et al.~\cite{19} present a contrasting case: on large-scale Chinese
legal case retrieval (LeCaRDv2), BM25 consistently outperforms
off-the-shelf dense models, indicating that dense dominance is
corpus- and language-dependent.

\subsection{Hybrid Retrieval}

Rayo et al.~\cite{3} show that combining BM25 with fine-tuned sentence
transformers substantially improves Recall@10 and MAP@10 on regulatory
corpora.
LegalRAG~\cite{10} demonstrates that hybrid systems outperform
single-mode systems for multilingual legal texts.
HyPA-RAG~\cite{2} introduces adaptive parameter adjustment for hybrid
retrieval in AI legal and policy applications.
SPLADE~\cite{5} and its efficiency-optimised variant~\cite{4} provide
sparse learned representations that address vocabulary mismatch while
maintaining inverted-index compatibility.
The original RRF algorithm~\cite{15} demonstrates that rank-based
fusion consistently outperforms score-based methods without requiring
score normalisation; Bruch et al.~\cite{18} further show that RRF
parameters are sensitive to corpus size and do not transfer reliably
from TREC-scale collections to small specialised domains.

\subsection{Retrieval Reliability and Legal RAG}

Reuter et al.~\cite{12} identify Document-Level Retrieval Mismatch
(DRM) as a systematic failure mode of dense retrievers on large legal
datasets, where shared structural boilerplate causes globally
irrelevant documents to rank highly.
Ho et al.~\cite{13} show that doctrinal accuracy requires structural
legal signals beyond semantic similarity.
CBR-RAG~\cite{6} and LexRAG~\cite{7} demonstrate that multi-turn
legal grounding requires retrieval robust to topic drift and implicit
definitional dependencies.

% ============================================================
\section{System Architecture}

\subsection{Corpus Construction}

The corpus is extracted from the EU CELLAR repository in two phases.
First, a SPARQL query retrieves all work URIs associated with
21~EuroVoc climate and energy concept identifiers.
Second, per-document manifestation queries resolve Formex~XML
(\texttt{fmx4}) URLs; an HTML fallback is used when no Formex
manifestation exists.

Formex~XML is parsed with \texttt{lxml}.
Each \texttt{<ARTICLE>} element becomes one corpus chunk; text is
reconstructed from \texttt{<ALINEA>} leaf nodes to avoid duplication.
Records are validated against a Pydantic schema; placeholders are
discarded.
Five articles exceeding the 8{,}192-token context window (all
Article~1 of amending instruments) are split at EU amendment-point
boundaries; sub-chunks inherit the parent \texttt{celex\_id}.

The final corpus contains \textbf{1{,}166} article-level provisions
from \textbf{72}~CELEX instruments covering ETS, CBAM, Taxonomy,
LULUCF, F-Gas, Maritime MRV, ESR, the European Climate Law, Energy
Efficiency, CSDDD, and the Governance Regulation.

\subsection{Sparse Retrieval --- BM25}

BM25Okapi (rank-bm25~0.2.2) with $k_1{=}1.5$, $b{=}0.75$.
A custom legal tokeniser preserves hyphenated statutory terms
(\textit{net-zero}, \textit{CBAM-certificate}) as single tokens.
No stemming; no stopword removal: EU definitions are terminologically
precise, and stopword removal breaks fixed legal phrases such as
\textit{free of charge} and \textit{not later than}.
At query time BM25 scores all 1{,}166 documents and returns the
top~100 ranked results.

\subsection{Dense Retrieval --- nomic-embed-text-v1.5}

\nomic is a pure bi-encoder~\cite{17} (768-dim, 8{,}192-token context window)
loaded via \texttt{sentence-transformers}.
It was selected because: (i)~it is fully open-weight; (ii)~its
8{,}192-token window accommodates EU articles without truncation for
all but five articles; (iii)~it is a \emph{pure} bi-encoder with no
sparse or late-interaction components, ensuring a clean
sparse-vs-dense comparison.
Models trained jointly for dense, sparse, and late-interaction
retrieval (e.g., BGE-M3) would contaminate this comparison because
their dense vectors already encode lexical signals.

Encoding uses asymmetric task prefixes required by the model:
\texttt{"search\_query:"} for queries and \texttt{"search\_document:"}
for indexed articles.
Corpus vectors are L2-normalised and stored in a FAISS
\texttt{IndexFlatIP}; inner product on normalised vectors is equivalent
to cosine similarity.
At query time, FAISS returns the top~100 exact nearest neighbours.

\subsection{Reciprocal Rank Fusion}

Ranked lists from BM25 and \nomic are merged via weighted Reciprocal
Rank Fusion.
For each candidate document~$d$:

\begin{equation}
  \mathrm{RRF}(d) = \frac{w_s}{k + \mathrm{rank}_s(d)}
                  + \frac{w_d}{k + \mathrm{rank}_d(d)}
  \label{eq:rrf}
\end{equation}

\noindent where $k{=}20$ is the smoothing constant,
$w_s{=}1$ (sparse weight), $w_d{=}5$ (dense weight), and
$\mathrm{rank}_r(d) = \infty$ if $d$ is absent from retriever~$r$.
Both retrievers are invoked concurrently via
\texttt{ThreadPool\-Executor(max\_workers=2)}.
The top~5 fused results are returned to the generation layer.

RRF is preferred over raw score fusion because BM25 scores are
unbounded positive reals while cosine similarities lie in
$[-1,1]$; combining them requires scale assumptions that may be
violated~\cite{15}.
The dense weight $w_d{=}5$ was tuned on a held-out validation split
of 18~queries; the test set of 53~queries was evaluated exactly once
after tuning.

% ============================================================
\section{Experimental Setup}

\subsection{Benchmark}

The benchmark contains 71~manually constructed queries covering all
eleven EU Climate Law instrument families.
Unlike broad multi-task legal benchmarks~\cite{20} that aggregate
hundreds of tasks across jurisdictions, this benchmark targets a single
tightly bounded corpus to enable precise evaluation of statute-level
retrieval.
Relevance is binary at the CELEX instrument level; a document is
relevant if it contains the provision that directly and materially
answers the query.
CELEX identifiers with corrigendum suffixes (e.g.,
\texttt{32003L0087R(02)}) are normalised by stripping the
\texttt{R(...)} suffix before matching, and duplicate CELEX hits within
a single result list are deduplicated to prevent NDCG inflation.

A stratified 70/30 split by query difficulty (seed~42) produces
18~validation queries (used for weight tuning) and a sealed
\textbf{53-query test set} used for all reported results.

\subsection{Metrics}

$\HR{5}$: fraction of queries with at least one relevant CELEX ID in
the top~5 results.
$\MRR{5}$: mean reciprocal rank of the first relevant result within
the top~5.
$\NDCG{5}$: normalised discounted cumulative gain at cut-off~5 with
binary relevance gains.
Wilson 95\,\% confidence intervals are reported for $\HR{5}$.
Pairwise statistical significance uses the Wilcoxon signed-rank test
on per-query MRR scores ($\alpha{=}0.05$, two-tailed, $n{=}53$).

\subsection{Hardware and Software}

All experiments run on an Apple~M1~Pro (CPU-only).
Python~3.14.4; \texttt{sentence-transformers}~5.5.1;
\texttt{faiss-cpu}~1.14.2; \texttt{rank-bm25}~0.2.2;
\texttt{pydantic}~2.13.4.
All library versions are pinned in \texttt{pyproject.toml}.

% ============================================================
\section{Results}

\subsection{Retrieval Effectiveness}

Table~\ref{tab:effectiveness} reports the main results on the
53-query test set.

\begin{table}[!t]
\renewcommand{\arraystretch}{1.3}
\caption{Retrieval effectiveness on the 53-query test set.
  Wilson 95\,\% CIs in parentheses for $\HR{5}$.
  $^\dagger p{=}0.0005$ vs.\ BM25;
  $^\ddagger p{=}0.004$ vs.\ BM25.}
\label{tab:effectiveness}
\centering
\small
\begin{tabular}{@{}lccc@{}}
\toprule
System & $\HR{5}$ & $\MRR{5}$ & $\NDCG{5}$ \\
\midrule
BM25
  & 0.698 (0.566--0.811)
  & 0.528
  & 0.494 \\
\nomic
  & \textbf{0.925} (0.849--0.981)$^\ddagger$
  & 0.743
  & 0.709 \\
Hybrid RRF
  & 0.906 (0.811--0.981)$^\dagger$
  & \textbf{0.773}
  & \textbf{0.721} \\
\bottomrule
\end{tabular}
\end{table}

Dense retrieval achieves the highest $\HR{5}$ (0.925 vs.\ 0.906 for
hybrid; vs.\ 0.698 for BM25).
The hybrid achieves the highest $\MRR{5}$ (0.773) and $\NDCG{5}$
(0.721), indicating superior ranking of relevant documents within the
top~5 even when coverage is marginally lower than dense alone.

All pairwise comparisons involving BM25 are statistically significant:
hybrid vs.\ BM25 ($p{=}0.0005$, Cohen's $d{=}0.579$, medium effect);
dense vs.\ BM25 ($p{=}0.004$, $d{=}0.433$, small effect).
The hybrid vs.\ dense comparison is not significant
($p{=}0.614$, $d{=}0.120$, negligible effect).

\subsection{Hard Failure Analysis}

Table~\ref{tab:reliability} breaks down query outcomes across systems.

\begin{table}[!t]
\renewcommand{\arraystretch}{1.3}
\caption{Per-query hit pattern breakdown ($n{=}53$).}
\label{tab:reliability}
\centering
\small
\begin{tabular}{@{}lcc@{}}
\toprule
Category & $n$ & Fraction \\
\midrule
All three systems hit        & 34 & 0.642 \\
Dense hits; BM25 misses      & 15 & 0.283 \\
Hybrid hits; dense misses    &  0 & 0.000 \\
Hybrid misses; dense hits    &  4 & 0.075 \\
All three miss               &  5 & 0.094 \\
\bottomrule
\end{tabular}
\end{table}

Dense alone recovers 15~queries that BM25 misses entirely, confirming
that semantic generalisation is the primary source of hybrid benefit.
The hybrid never strictly improves over dense (0~queries where hybrid
hits and dense misses), while it regresses on 4~queries where BM25
lexical errors displace a dense rank-1 result.
Five queries are hard failures for all three systems, likely requiring
provisions from instruments outside the 72-instrument corpus.

\subsection{Ranking Quality}

The hybrid's advantage over dense on $\MRR{5}$ ($+0.031$) and
$\NDCG{5}$ ($+0.013$) indicates that when both systems hit within
the top~5, the hybrid tends to rank the relevant document higher.
This is consistent with the complementarity hypothesis: BM25's
exact-term matching corrects cases where \nomic places the correct
instrument at rank~2 or~3 behind semantically similar alternatives.

% ============================================================
\section{Discussion}

Hybrid RRF significantly outperforms BM25 ($p{=}0.0005$, medium
effect), confirming that semantic retrieval is essential for EU Climate
Law queries where user vocabulary diverges from statutory text.
However, the hybrid does not significantly outperform dense alone
($p{=}0.614$): the asymmetric weight ($w_d{=}5$) prevents BM25 errors
from degrading the dense ranking, but BM25 contributes no unique hits
that dense misses.

The $9.4\,\%$ hard-failure rate is primarily driven by queries
requiring instruments outside the current corpus rather than retrieval
strategy limitations.
The efficiency overhead of the hybrid is negligible: concurrent
retrieval via \texttt{ThreadPoolExecutor} means hybrid latency is
dominated by the dense component (${\approx}35$\,ms median on CPU).
The combined index occupies only 9.7\,MB on disk.

The literature default of $k{=}60$ with equal weights was found in
preliminary experiments to actively degrade hybrid performance
($\HR{5}{=}0.887$) because BM25 lexical errors displaced dense rank-1
results.
Reducing $k$ to~20 and weighting dense at~5 protects the dense
ranking from BM25 interference.
This corpus-size dependence confirms that RRF parameters calibrated
on TREC-scale collections do not transfer to small specialised
corpora~\cite{18}.

% ============================================================
\section{Conclusion}

Dense retrieval (\nomic) significantly outperforms BM25
($p{=}0.004$, $\HR{5}{=}0.925$ vs.\ $0.698$).
Weighted hybrid RRF achieves the best ranking metrics ($\MRR{5}{=}0.773$,
$\NDCG{5}{=}0.721$) and significantly outperforms BM25
($p{=}0.0005$, medium effect), but does not significantly outperform
dense alone ($p{=}0.614$).

The primary practical conclusion is that \nomic dense retrieval is
already near-optimal for first-stage statute retrieval in this domain,
and that hybrid fusion provides a statistically non-significant but
directionally positive improvement in ranking quality.
For legal RAG systems where hard failures are the principal risk,
corpus expansion is likely to provide greater gains than retrieval
strategy change.

Future work will: (i)~compute Cohen's~$\kappa$ from independent
relevance annotation; (ii)~evaluate a cross-encoder re-ranking stage
on the top-5 hybrid output; (iii)~extend the corpus to consolidated
instrument versions to eliminate amendment-article boilerplate; and
(iv)~explore multilingual retrieval exploiting \nomic's long-context
capability.

\section*{Generative AI Disclosure}

Generative AI tools were used solely for code scaffolding and to
assist with clarity of writing.
All research ideas, experimental designs, relevance judgements, and
reported results are entirely our own.

\begin{thebibliography}{20}

\bibitem{1}
L.~Mori, C.~Sousa de Oliveira, Y.~Yih, and M.~Ventresca,
``Assessing the performance gap between lexical and semantic models
for information retrieval with formulaic legal language,''
in \emph{Proc.\ 20th Int.\ Conf.\ Artificial Intelligence and Law},
pp.~114--128, 2025.

\bibitem{2}
R.~Kalra et al.,
``HyPA-RAG: A hybrid parameter adaptive retrieval-augmented generation
system for AI legal and policy applications,''
in \emph{Proc.\ Workshop on Customizable NLP (CustomNLP4U)},
pp.~237--256, 2024.

\bibitem{3}
J.~Rayo, R.~de La Rosa, and M.~Garrido,
``A hybrid approach to information retrieval and answer generation for
regulatory texts,''
\emph{RegNLP Workshop}, 2025.

\bibitem{4}
C.~Lassance and S.~Clinchant,
``An efficiency study for SPLADE models,''
in \emph{Proc.\ 45th Int.\ ACM SIGIR},
pp.~2220--2226, 2022.

\bibitem{5}
T.~Formal, C.~Lassance, B.~Piwowarski, and S.~Clinchant,
``SPLADE v2: Sparse lexical and expansion model for information
retrieval,'' \emph{arXiv:2109.10086}, 2021.

\bibitem{6}
N.~Wiratunga et al.,
``CBR-RAG: Case-based reasoning for retrieval augmented generation in
LLMs for legal question answering,''
in \emph{Int.\ Conf.\ Case-Based Reasoning}, pp.~445--460, 2024.

\bibitem{7}
H.~Li et al.,
``LexRAG: Benchmarking retrieval-augmented generation in multi-turn
legal consultation conversation,''
in \emph{Proc.\ 48th Int.\ ACM SIGIR}, pp.~3606--3615, 2025.

\bibitem{8}
L.~Zheng et al.,
``A reasoning-focused legal retrieval benchmark,''
in \emph{Proc.\ 2025 Symp.\ Computer Science and Law},
pp.~169--193, 2025.

\bibitem{9}
A.~Chouhan and M.~Gertz,
``LexDrafter: Terminology drafting for legislative documents using
retrieval augmented generation,''
in \emph{Proc.\ LREC-COLING 2024}, pp.~10448--10458, 2024.

\bibitem{10}
M.~R.~Kabir et al.,
``LegalRAG: A hybrid RAG system for multilingual legal information
retrieval,''
in \emph{Proc.\ IJCNN 2025}, IEEE, 2025.

\bibitem{11}
H.~T.~Nguyen et al.,
``NOWJ@COLIEE 2025: A multi-stage framework integrating embedding
models and LLMs for legal retrieval and entailment,''
\emph{arXiv:2509.08025}, 2025.

\bibitem{12}
M.~Reuter et al.,
``Towards reliable retrieval in RAG systems for large legal datasets,''
in \emph{Proc.\ Natural Legal Language Processing Workshop 2025},
pp.~17--30, 2025.

\bibitem{13}
J.~Ho, A.~Colby, and W.~Fisher,
``Incorporating legal structure in retrieval-augmented generation:
A case study on copyright fair use,''
\emph{arXiv:2505.02164}, 2025.

\bibitem{14}
M.~Park, H.~Oh, E.~Choi, and W.~Hwang,
``LRAGE: Legal retrieval augmented generation evaluation tool,''
\emph{arXiv:2504.01840}, 2025.

\bibitem{15}
G.~V.~Cormack, C.~L.~A.~Clarke, and S.~B\"uttcher,
``Reciprocal rank fusion outperforms condorcet and individual rank
learning methods,''
in \emph{Proc.\ 32nd Int.\ ACM SIGIR}, pp.~758--759, 2009.
\url{https://doi.org/10.1145/1571941.1572114}

\bibitem{16}
N.~Thakur, N.~Reimers, A.~R\"uckl\'e, A.~Srivastava, and I.~Gurevych,
``BEIR: A heterogeneous benchmark for zero-shot evaluation of
information retrieval models,''
in \emph{Advances in Neural Information Processing Systems}, vol.~34,
pp.~17398--17411, 2021.
\url{https://arxiv.org/abs/2104.08663}

\bibitem{17}
V.~Karpukhin et al.,
``Dense passage retrieval for open-domain question answering,''
in \emph{Proc.\ 2020 Conf.\ Empirical Methods in Natural Language
Processing (EMNLP)}, pp.~6769--6781, 2020.
\url{https://doi.org/10.18653/v1/2020.emnlp-main.550}

\bibitem{18}
S.~Bruch, S.~Gai, and A.~Ingber,
``An analysis of fusion functions for hybrid retrieval,''
\emph{ACM Trans.\ Inf.\ Syst.}, vol.~42, no.~1, pp.~20:1--20:35, 2024.
\url{https://doi.org/10.1145/3596512}

\bibitem{19}
H.~Li et al.,
``LeCaRDv2: A large-scale Chinese legal case retrieval dataset,''
in \emph{Proc.\ 47th Int.\ ACM SIGIR}, pp.~2251--2260, 2024.
\url{https://doi.org/10.1145/3626772.3657887}

\bibitem{20}
N.~Guha et al.,
``LegalBench: A collaboratively built benchmark for measuring legal
reasoning in large language models,''
in \emph{Advances in Neural Information Processing Systems}, vol.~36,
2023.
\url{https://arxiv.org/abs/2308.11462}

\end{thebibliography}

\end{document}
